\textcolor{blue}{3. [10 min/15 puntos] Explica en un párrafo las diferencias entre el cómputo concurrente, paralelo
y distribuido.}\\

El cómputo paralelo consiste en realizar varias tareas que se ejecutan simultáneamente, normalmente utilizando varios núcleos o procesadores. En cambio, en el cómputo distribuido, las tareas se ejecutan en diferentes computadoras o nodos conectados mediante una red, los cuales trabajan en conjunto para resolver un problema. Por otro lado, el cómputo concurrente consiste en realizar y gestionar varias tareas que comparten recursos dentro de un mismo sistema, permitiendo que estas avancen de manera simultánea o intercalada. Por lo tanto, podemos considerar que el cómputo paralelo y el distribuido son formas de cómputo concurrente, ya que ambos implican la ejecución o gestión de múltiples tareas que pueden avanzar de manera independiente.\\
